\documentclass{article}

% Packages for math and theorems
\usepackage{amsmath}
\usepackage{amsthm}
\usepackage{amssymb}
\usepackage{graphicx}
\usepackage{float}
\usepackage{geometry}
\geometry{margin=1.5in}

% Theorem style
\newtheorem{theorem}{Theorem}
\newtheorem{lemma}{Lemma}
\newtheorem{corollary}{Corollary}

% Package for drawings (Asymptote for vector graphics)
\usepackage{asymptote}

% Title and author
\title{Geometry Theorem Template}
\author{Hanson Jiang}
\date{\today}

\begin{document}

\maketitle

\section{Prop. 4}

\subsection{Theorem Statement}

\begin{theorem}[Prop. 4: SAS Congruence]
  If ABC and DEF are two triangles such that $\overline{AB} \cong
  \overline{DE}$, $\overline{AC} \cong \overline{DF}$, and
  $\angle$ BAC $\cong$ $\angle$ EDF, then $\triangle$ ABC $\cong$
  $\triangle$ DEF.
\end{theorem}

\subsection{Diagram}

\begin{figure}[h]
  \centering
  \begin{asy}
    size(8cm);

    // Triangle 1 (ABC)
    pair A = (0, 0);
    pair B = (3, 0);
    pair C = (1, 2.5);

    draw(A--B--C--cycle);

    // Triangle 2 (DEF)
    pair D = (5, 0);
    pair E = (8, 0);
    pair F = (6, 2.5);

    draw(D--E--F--cycle);

    // Labels
    label("$A$", A, SW);
    label("$B$", B, SE);
    label("$C$", C, N);
    label("$D$", D, SW);
    label("$E$", E, SE);
    label("$F$", F, N);
  \end{asy}
  \caption{SAS Congruence}
  \label{fig:theorem-diagram}
\end{figure}

\section{Prop. 5}

\subsection{Theorem Statement}

\begin{theorem}[Prop. 5: Isosceles Triangle Theorem Variant]
  If $\overline{AB} \cong \overline{AC}$, then $\angle$ DBC $\cong$
  $\angle$ ECB.
\end{theorem}

\subsection{Diagram}

\begin{figure}[H]
  \centering
  \begin{asy}
    size(5cm);

    // Isosceles triangle (two equal sides) - tall
    pair A = (1, 3);      // Left base vertex
    pair B = (0, 0);      // Right base vertex
    pair C = (2, 0);      // Top vertex (apex) - makes AC = BC

    // Draw the triangle
    draw(A--B--C--cycle);

    // Extend the equal sides slightly beyond the base
    pair D = B + 0.3 * (B - A);  // Extend BA past B
    pair E = C + 0.3 * (C - A);  // Extend CA past C

    draw(B--D, black);  // Extended side BA
    draw(C--E, black);  // Extended side CA

    // Labels
    label("A", A, N);
    label("B", B, W);
    label("C", C, E);
    label("D", D, SW);
    label("E", E, SE);
  \end{asy}
  \caption{Variant of the Isosceles Triangle Theorem}
  \label{fig:theorem-diagram}
\end{figure}

\subsection{Corollaries or Additional Notes}

\begin{corollary}
  Euclid does not prove the version of the isosceles triangle theorem
  that is normally taught today.
\end{corollary}

\section{Prop. 6}

\subsection{Theorem Statement}

\begin{theorem}[Prop. 6: Converse of the Isosceles Triangle Theorem]
  If $\triangle$ ABC is a triangle such that $\angle$ ABC $\cong$
  $\angle$ ACB, then $\overline{AB} \cong \overline{AC}$.
\end{theorem}

\subsection{Diagram}

\begin{figure}[h]
  \centering
  \begin{asy}
    size(5cm);

    // Isosceles triangle (two equal sides) - tall
    pair A = (1, 3);      // Left base vertex
    pair B = (0, 0);      // Right base vertex
    pair C = (2, 0);      // Top vertex (apex) - makes AC = BC

    // Draw the triangle
    draw(A--B--C--cycle);

    // Labels
    label("A", A, N);
    label("B", B, SW);
    label("C", C, SE);

  \end{asy}
  \caption{Converse of the Isosceles Triangle Theorem}
  \label{fig:theorem-diagram}
\end{figure}

\subsection{Corollaries or Additional Notes}

\begin{corollary}
  Any additional corollaries or notes can go here.
\end{corollary}

\end{document}